\chapter{Antecedentes}

\section{Marco Teórico}

\subsection{Computación en la nube (Cloud Computing)}

\subsubsection{Definición de computación en la nube}

La computación en la nube se define como un modelo que permite el acceso ubicuo, conveniente y bajo demanda, mediante red, a un conjunto compartido de recursos computacionales configurables que pueden aprovisionarse y liberarse rápidamente con un mínimo esfuerzo administrativo o interacción con el proveedor del servicio \parencite{Mell2011}.

\paragraph{Características principales}

De acuerdo con MELL y GRANCE (2011), la computación en la nube posee cinco características esenciales que permiten diferenciar este modelo de otros enfoques tradicionales de provisión de recursos informáticos. En primer lugar, el autoservicio bajo demanda permite que los usuarios aprovisionen recursos de manera automática según sus necesidades, sin requerir intervención directa del proveedor. En segundo lugar, el acceso amplio a la red posibilita utilizar los servicios desde distintos dispositivos conectados a Internet. A su vez, la agrupación de recursos permite compartir infraestructura física entre múltiples usuarios mediante mecanismos de virtualización, optimizando el uso de los recursos disponibles. Otra característica es la elasticidad rápida, que permite aumentar o disminuir dinámicamente la capacidad disponible en función de la demanda. Finalmente, el servicio medido habilita el monitoreo y control del consumo de recursos, permitiendo modelos de facturación basados en el uso efectivo del servicio. Estas características contribuyen a mejorar la eficiencia operativa y la flexibilidad en la gestión de infraestructura tecnológica.

\subsection{Modelos de servicio}

Según MELL y GRANCE (2011), la computación en la nube se estructura en tres modelos principales de servicio: Software as a Service (SaaS), donde el usuario utiliza aplicaciones alojadas por el proveedor; Platform as a Service (PaaS), que permite desarrollar y desplegar aplicaciones utilizando una plataforma administrada; e Infrastructure as a Service (IaaS), que ofrece acceso a recursos de infraestructura como procesamiento, almacenamiento y redes.

\subsection{Modelos de despliegue}

Según MELL y GRANCE (2011), la computación en la nube puede implementarse mediante distintos modelos de despliegue: nube privada (private cloud), destinada al uso exclusivo de una organización; nube comunitaria (community cloud), compartida por organizaciones con intereses o requisitos comunes; nube pública (public cloud), disponible para uso abierto del público general; y nube híbrida (hybrid cloud), que combina dos o más infraestructuras de nube manteniendo interoperabilidad entre ellas.

En síntesis, la computación en la nube constituye un paradigma orientado a la provisión flexible de recursos tecnológicos mediante Internet. Sus características, modelos de servicio y formas de despliegue permiten adaptar la infraestructura a distintos contextos organizacionales \parencite{Mell2011}.

\subsection{Arquitectura y administración de infraestructura cloud}

\subsubsection{Concepto de infraestructura cloud}

La infraestructura cloud se compone de un conjunto de recursos computacionales configurables y compartidos, entre los que se incluyen capacidades de procesamiento, almacenamiento, redes y otros servicios tecnológicos accesibles mediante Internet. Estos recursos pueden aprovisionarse y liberarse rápidamente de acuerdo con la demanda, permitiendo una mayor flexibilidad y escalabilidad respecto de modelos tradicionales basados en infraestructura física local \parencite{Mell2011}.

En entornos empresariales, la provisión de infraestructura cloud suele ser ofrecida por plataformas especializadas que permiten desplegar y administrar recursos tecnológicos bajo demanda, reduciendo la necesidad de mantener infraestructura propia. Entre los principales proveedores de servicios cloud se encuentran plataformas comerciales ampliamente utilizadas en el mercado, como Amazon Web Services (AWS), Microsoft Azure, Google Cloud Platform (GCP) y Huawei Cloud.

Amazon Web Services define la computación en la nube como la entrega bajo demanda de recursos de TI mediante Internet con un modelo de pago por consumo \parencite{AWS2024}.

\subsubsection{Administración de infraestructura cloud}

La administración de infraestructura cloud comprende el conjunto de actividades orientadas al aprovisionamiento, configuración, supervisión y optimización de recursos computacionales ofrecidos por plataformas de servicios en la nube. A través de mecanismos de administración centralizada, automatización y escalabilidad dinámica, los proveedores cloud permiten crear, modificar y eliminar recursos de infraestructura según las necesidades operativas de cada organización, reduciendo la dependencia del mantenimiento directo del hardware físico \parencite{Mell2011, AWS2023}.

Entre las principales actividades involucradas en la administración de infraestructura cloud se encuentran la gestión de servidores virtuales, la configuración de redes y sistemas de almacenamiento, el monitoreo del consumo de recursos, la escalabilidad de servicios en función de la demanda, el control de costos operativos y la automatización de despliegues \parencite{AWS2026}. Estas prácticas permiten mejorar la disponibilidad y el rendimiento de los servicios, optimizar el uso de los recursos tecnológicos y favorecer una mayor flexibilidad operativa dentro de los entornos cloud.

\subsection{Escalabilidad y elasticidad}

La capacidad de adaptación frente a cambios en la demanda constituye una de las principales ventajas de las arquitecturas basadas en computación en la nube. En este contexto, la escalabilidad y la elasticidad permiten ajustar el consumo de recursos y mantener niveles adecuados de rendimiento, disponibilidad y eficiencia operativa.

La escalabilidad hace referencia a la capacidad de un sistema para incrementar o reducir los recursos disponibles con el objetivo de responder a variaciones en la carga de trabajo manteniendo niveles aceptables de desempeño. En entornos cloud, esta capacidad puede implementarse mediante mecanismos de escalado vertical, aumentando la capacidad de recursos existentes, o mediante escalado horizontal, incorporando nuevas instancias o componentes dentro de la infraestructura. Esta propiedad permite que los sistemas puedan adaptarse al crecimiento o disminución de la demanda sin comprometer su funcionamiento \parencite{Ullah2018}.

Por otra parte, la elasticidad corresponde a la capacidad de aprovisionar y liberar recursos de manera rápida y adaptable frente a cambios en la demanda, generando para el consumidor la percepción de contar con recursos prácticamente ilimitados disponibles cuando sean requeridos \parencite{Mell2011}.

Si bien ambos conceptos se encuentran estrechamente relacionados, presentan diferencias conceptuales. Mientras que la escalabilidad se orienta principalmente al crecimiento o reducción de capacidad para soportar variaciones en la carga del sistema, la elasticidad incorpora además la capacidad de realizar dichos ajustes de forma dinámica y en correspondencia temporal con la demanda efectiva, buscando optimizar el uso de recursos disponibles \parencite{Ullah2018, Mell2011}.

\subsection{Monitoreo}

\subsubsection{Concepto de monitoreo}

El monitoreo comprende el conjunto de actividades orientadas a recopilar, procesar y analizar información sobre el estado y comportamiento de sistemas e infraestructura con el objetivo de evaluar su rendimiento, disponibilidad y funcionamiento operativo. En entornos de computación en la nube, el monitoreo adquiere mayor relevancia debido a características como la escalabilidad, la elasticidad y el aprovisionamiento dinámico de recursos, requiriendo mecanismos capaces de obtener información continua para facilitar la administración y toma de decisiones sobre la infraestructura \parencite{Ward2014}.

El monitoreo de infraestructura cloud comprende el conjunto de procesos orientados a recopilar, procesar y analizar información sobre el estado y comportamiento de los recursos desplegados en entornos de computación en la nube con el objetivo de evaluar su rendimiento, disponibilidad y funcionamiento operativo. De acuerdo con Ward y Barker (2014), el proceso de monitoreo se compone principalmente de tres etapas: la recolección de datos, el análisis de la información obtenida y la toma de decisiones basada en los resultados del análisis. Para llevar adelante estas actividades, las plataformas de monitoreo emplean mecanismos como agentes distribuidos, servidores de monitoreo, sistemas de recolección mediante consultas o envío de métricas, almacenamiento de estados del sistema, generación de alertas y herramientas de visualización. También pueden utilizar métricas de rendimiento, utilización de recursos, disponibilidad, registros de eventos y otros indicadores operativos para facilitar la administración y el control de entornos cloud \parencite{Ward2014}.

El monitoreo en entornos cloud puede realizarse sobre diferentes niveles de la infraestructura y los servicios desplegados. De acuerdo con Ward y Barker (2014), las estrategias de monitoreo suelen abarcar componentes físicos y virtuales, recursos computacionales, redes, almacenamiento, aplicaciones y servicios ofrecidos al usuario final. Esta organización multinivel permite obtener una visión más completa del comportamiento del entorno y facilita la identificación de problemas de rendimiento, disponibilidad o utilización de recursos dentro de arquitecturas distribuidas \parencite{Ward2014}.

En consecuencia, el monitoreo adquiere un rol central dentro de la administración de infraestructura cloud, ya que permite obtener visibilidad sobre el comportamiento de recursos y servicios en entornos caracterizados por cambios dinámicos y distribución de componentes. Debido a la complejidad y variabilidad de estos entornos, las estrategias de monitoreo deben combinar distintos mecanismos de recolección y análisis de información para adaptarse a los objetivos operativos y facilitar la gestión eficiente de la infraestructura. En este contexto, el monitoreo constituye una base fundamental para tareas posteriores de optimización, automatización y mejora continua de los sistemas desplegados en la nube \parencite{Ward2014}.

\subsection{Observabilidad}

La observabilidad corresponde a la capacidad de comprender el estado interno y el comportamiento de un sistema a partir del análisis de las señales que éste genera externamente. En entornos cloud y arquitecturas distribuidas, la observabilidad amplía el alcance del monitoreo tradicional al no limitarse únicamente a detectar eventos previamente definidos, sino que busca proporcionar información suficiente para investigar comportamientos inesperados y comprender el origen de incidentes. Para ello, las plataformas de observabilidad suelen apoyarse en la recopilación y correlación de tres fuentes principales de información: métricas, que permiten representar el comportamiento del sistema mediante indicadores cuantitativos; registros de eventos (logs), que aportan detalle sobre sucesos ocurridos; y trazabilidad distribuida (distributed tracing), que permite reconstruir el recorrido de solicitudes entre múltiples componentes del sistema \parencite{Kratzke2022}.

En este sentido, mientras que el monitoreo se orienta principalmente al seguimiento continuo del estado de la infraestructura y la generación de alertas, la observabilidad busca proporcionar el nivel de información necesario para interpretar el comportamiento del sistema y facilitar el diagnóstico de problemas complejos \parencite{Kratzke2022}.

\subsection{Gestión y optimización de costos cloud}

La gestión y optimización de costos en entornos de computación en la nube comprende el conjunto de prácticas orientadas a supervisar, analizar y controlar el consumo de recursos tecnológicos con el objetivo de maximizar el valor generado por la infraestructura utilizada. Debido al modelo de consumo bajo demanda característico de la nube, las organizaciones requieren mecanismos que permitan mantener visibilidad sobre el uso de recursos, comprender el impacto económico de las decisiones técnicas y ajustar el gasto en función de objetivos operativos y de negocio \parencite{FinOps2026}.

En este contexto surge FinOps, entendido como un marco operativo y una práctica organizacional orientada a maximizar el valor del uso de tecnologías mediante la colaboración entre equipos técnicos, financieros y de negocio. A diferencia de enfoques tradicionales centrados únicamente en la reducción de costos, FinOps propone generar responsabilidad financiera compartida y favorecer la toma de decisiones basada en datos para equilibrar variables como velocidad, calidad y gasto operativo \parencite{FinOps2026}.

De acuerdo con el FinOps Framework, la gestión financiera de la nube se desarrolla como un proceso iterativo compuesto por tres etapas principales: Inform, orientada a generar visibilidad y comprensión del consumo y los costos; Optimize, enfocada en identificar oportunidades de eficiencia y reducción del desperdicio; y Operate, destinada a incorporar métricas, gobierno y procesos continuos para sostener la optimización en el tiempo. Entre las prácticas asociadas se incluyen la asignación de costos, presupuestación, análisis del consumo, definición de indicadores y establecimiento de mecanismos de control y seguimiento del gasto tecnológico \parencite{FinOpsFramework2026}.

La optimización de recursos cloud comprende el conjunto de prácticas orientadas a asignar, ajustar y administrar los recursos computacionales de manera eficiente con el objetivo de maximizar el rendimiento, mantener niveles adecuados de disponibilidad y reducir el desperdicio de capacidad y costos operativos. En entornos cloud, este proceso busca alinear el aprovisionamiento de infraestructura con las necesidades reales de las cargas de trabajo mediante mecanismos como el ajuste de capacidad (rightsizing), el escalado automático, la redistribución de cargas y el monitoreo continuo del consumo de recursos. Asimismo, la optimización requiere visibilidad permanente sobre el comportamiento del entorno para identificar recursos subutilizados o sobreaprovisionados y favorecer decisiones que permitan equilibrar desempeño, eficiencia y costo dentro de la infraestructura \parencite{IBM2026, FinOpsFramework2026}.

La optimización de recursos complementa las prácticas de gestión y optimización de costos, ya que no se orienta únicamente a disminuir el gasto económico, sino también a asegurar que los recursos consumidos generen valor operativo y se encuentren alineados con los objetivos del negocio \parencite{IBM2026, FinOps2026}.

\subsection{Infrastructure as Code (IaC)}

\subsubsection{Concepto de Infrastructure as Code}

La Infraestructura como Código (Infrastructure as Code -- IaC) corresponde a una práctica orientada a definir, aprovisionar y administrar infraestructura tecnológica mediante archivos de configuración y procesos automatizados en lugar de realizar configuraciones manuales sobre los recursos. Bajo este enfoque, la infraestructura es tratada de forma similar al código fuente de una aplicación, permitiendo aplicar mecanismos como control de versiones, reutilización, pruebas y despliegues repetibles para mejorar la consistencia y reducir errores operativos. Asimismo, IaC busca evitar diferencias entre entornos y favorecer una administración más eficiente de recursos en arquitecturas modernas y entornos cloud \parencite{Microsoft2025}.

\subsubsection{Terraform}

Entre las herramientas utilizadas para implementar este enfoque se encuentra Terraform, una herramienta de Infraestructura como Código desarrollada por HashiCorp que permite definir y administrar recursos mediante archivos declarativos reutilizables. Terraform posibilita automatizar el aprovisionamiento y administración de infraestructura tanto en entornos locales como en múltiples proveedores cloud utilizando un flujo consistente para desplegar y mantener recursos a lo largo de su ciclo de vida \parencite{HashiCorp2026}.

En este sentido, IaC contribuye a incrementar la estandarización y reproducibilidad de la infraestructura, facilitando procesos de automatización y reduciendo la dependencia de configuraciones manuales dentro de entornos de computación en la nube \parencite{Microsoft2025, HashiCorp2026}.

\subsection{Machine Learning aplicado a infraestructura cloud}

\subsubsection{Concepto de Machine Learning}

El Machine Learning (aprendizaje automático) corresponde a una rama de la inteligencia artificial orientada al desarrollo de modelos capaces de identificar patrones y generar predicciones o decisiones a partir del análisis de datos, sin requerir que todas las reglas de comportamiento sean definidas explícitamente mediante programación tradicional. A través del entrenamiento sobre conjuntos de datos, los algoritmos de aprendizaje automático ajustan sus parámetros internos con el objetivo de mejorar su desempeño frente a una determinada tarea \parencite{Forghani2020}.

\subsubsection{Aprendizaje supervisado y no supervisado}

Dentro de las principales categorías de Machine Learning se encuentran el aprendizaje supervisado y el aprendizaje no supervisado. El aprendizaje supervisado se basa en el uso de conjuntos de datos etiquetados, donde cada dato de entrada se encuentra asociado a un resultado esperado, permitiendo que el modelo aprenda relaciones entre variables para realizar tareas de clasificación o predicción. Por el contrario, el aprendizaje no supervisado trabaja con datos sin etiquetar y busca descubrir patrones, agrupamientos o estructuras presentes en los datos sin contar previamente con resultados conocidos \parencite{Forghani2020}.

\subsubsection{Series temporales}

Las series temporales consisten en secuencias de observaciones registradas en distintos instantes y ordenadas cronológicamente. Habitualmente, estas observaciones se recolectan en intervalos de tiempo regulares, lo que permite analizar la evolución de una variable a lo largo del tiempo. Además, una serie temporal puede ser univariada cuando contiene una única variable por instante, o multivariada cuando registra múltiples variables en cada momento \parencite{Haben2023}.

En el contexto del monitoreo de infraestructura cloud, las métricas de utilización de recursos pueden representarse mediante series temporales, ya que valores como el uso de CPU, memoria, almacenamiento y tráfico de red son recolectados periódicamente. En InfraEye, estas métricas son obtenidas en intervalos de cinco minutos, permitiendo construir registros temporales sobre el comportamiento de los recursos. El análisis de estos datos permite obtener características estadísticas sobre diferentes ventanas temporales y utilizarlas posteriormente como entrada para el proceso de detección de anomalías.

El análisis mediante ventanas temporales permite resumir el comportamiento de las métricas durante períodos determinados. En lugar de considerar únicamente valores individuales, pueden calcularse medidas estadísticas como la media, el máximo, el mínimo, la desviación estándar y los percentiles de una ventana de observaciones. De esta manera, es posible caracterizar el comportamiento de utilización de un recurso durante un período y reducir la influencia de variaciones puntuales.

\subsection{Algoritmo Isolation Forest}

Isolation Forest es un algoritmo de aprendizaje automático no supervisado orientado a la detección de anomalías mediante el aislamiento de observaciones que presentan comportamientos diferentes respecto del conjunto de datos analizado. A diferencia de enfoques tradicionales basados en medidas de distancia o densidad, Isolation Forest identifica anomalías construyendo múltiples árboles de decisión aleatorios (Isolation Trees) que separan progresivamente los datos mediante particiones sucesivas. Bajo este enfoque, las observaciones anómalas tienden a requerir una menor cantidad de divisiones para quedar aisladas respecto de los datos considerados normales, permitiendo asignar una puntuación de anomalía a cada observación \parencite{Chen2023}.

Entre las principales ventajas de Isolation Forest se encuentran su baja complejidad computacional, la capacidad de operar eficientemente sobre conjuntos de datos de gran volumen y alta dimensionalidad, y el hecho de no requerir datos previamente etiquetados para su entrenamiento. Estas características favorecen su aplicación en escenarios donde los eventos anómalos son poco frecuentes o difíciles de identificar previamente. Además, al no depender directamente de cálculos de distancia entre observaciones, el algoritmo puede reducir el costo computacional respecto de otros enfoques tradicionales de detección de anomalías \parencite{Chen2023}.

En el contexto de INFRAEYE, Isolation Forest se aplica específicamente sobre las métricas de utilización de recursos recolectadas periódicamente desde la infraestructura cloud del usuario. Estas métricas incluyen la utilización de CPU, memoria, disco, tráfico de red entrante y tráfico de red saliente, las cuales son obtenidas en intervalos regulares de cinco minutos, conformando series temporales multivariadas. Dado que Isolation Forest no opera directamente sobre series temporales, se realiza una transformación que representa el comportamiento de cada instancia mediante características estadísticas calculadas sobre ventanas temporales: la media, el percentil 5 y el percentil 95 de cada métrica. De esta forma, cada instancia queda representada por 15 variables que describen tanto su nivel habitual de utilización como los niveles extremos que alcanza, permitiendo que el algoritmo identifique instancias cuyo comportamiento se desvía significativamente respecto del patrón de referencia \parencite{Chen2023}.

\subsection{Interfaces gráficas y visualización de datos}

La visualización de datos comprende el proceso de representar información mediante elementos gráficos con el objetivo de facilitar su interpretación, exploración y análisis por parte de los usuarios. A través de mecanismos como gráficos, paneles de control (dashboards), indicadores visuales y representaciones interactivas, la visualización permite transformar grandes volúmenes de datos en información más comprensible y útil para la toma de decisiones \parencite{Munzner2014}.

En entornos de computación en la nube, la visualización adquiere especial relevancia debido a la cantidad de información generada por métricas de consumo, monitoreo, utilización de recursos y costos operativos. La representación visual de estos datos permite identificar patrones de uso, detectar comportamientos anómalos y comprender el impacto de las decisiones sobre la infraestructura, favoreciendo procesos de optimización y gestión del entorno tecnológico \parencite{Munzner2014}.

Por otra parte, la experiencia del usuario (User Experience -- UX) aplicada a la visualización busca diseñar mecanismos que permitan a los usuarios interpretar información de manera eficiente y reducir la carga cognitiva asociada al análisis de datos. En este contexto, el uso adecuado de elementos visuales, jerarquías de información, interacción y organización de contenidos puede mejorar la velocidad de comprensión y favorecer una toma de decisiones más efectiva sobre sistemas complejos \parencite{Shneiderman1996}.

\subsection{DevOps y Automatización}

\subsubsection{Concepto de DevOps}

DevOps corresponde a un enfoque organizacional y técnico orientado a integrar las actividades de desarrollo de software (Development) y operaciones de infraestructura (Operations) con el objetivo de mejorar la velocidad, confiabilidad y calidad en la entrega de soluciones tecnológicas. Este enfoque promueve la colaboración entre equipos, la automatización de procesos y la adopción de prácticas que permitan reducir el tiempo entre el desarrollo y la puesta en producción de cambios, manteniendo la estabilidad y disponibilidad de los servicios \parencite{GoogleCloud2026}.

Dentro de este contexto, el profesional DevOps participa en el diseño, implementación y mantenimiento de procesos que permitan automatizar y optimizar el ciclo de vida del software y la infraestructura. Entre sus responsabilidades suelen incluirse la construcción y administración de pipelines de integración y entrega continua (CI/CD), la administración de infraestructura cloud, la automatización del aprovisionamiento mediante Infraestructura como Código (IaC), el monitoreo de servicios, la gestión de despliegues y el análisis de incidentes operativos \parencite{GoogleCloud2026}.

Las investigaciones del programa DevOps Research and Assessment (DORA) indican que la adopción de capacidades técnicas y organizacionales asociadas a DevOps contribuye a mejorar el rendimiento de entrega de software y el desempeño general de las organizaciones. Entre dichas capacidades se incluyen la automatización, la integración continua, la entrega continua, la administración eficiente de infraestructura cloud y el monitoreo continuo de los servicios \parencite{GoogleCloud2026}.

\subsubsection{Automatización de procesos}

La automatización corresponde al uso de procesos, herramientas y mecanismos tecnológicos destinados a ejecutar tareas de forma automática con mínima o nula intervención manual. En entornos de infraestructura y computación en la nube, la automatización busca reducir actividades repetitivas y operativas mediante la definición de reglas, configuraciones y flujos de trabajo capaces de aprovisionar, administrar y mantener recursos tecnológicos de manera consistente y reproducible \parencite{RedHat2026}.

Dentro de la administración de infraestructura cloud, la automatización puede aplicarse sobre actividades como el aprovisionamiento de recursos, despliegue de aplicaciones, configuración de entornos, monitoreo, recuperación ante fallos, escalabilidad y optimización del consumo de recursos. Este enfoque contribuye a disminuir errores derivados de configuraciones manuales, mejorar la velocidad de entrega de cambios y aumentar la capacidad de respuesta frente a variaciones en la demanda del sistema \parencite{RedHat2026}.

La automatización constituye uno de los principios fundamentales dentro de las prácticas DevOps, permitiendo integrar procesos de desarrollo y operaciones mediante mecanismos como integración continua (Continuous Integration -- CI), entrega continua (Continuous Delivery -- CD), Infraestructura como Código (Infrastructure as Code -- IaC) y despliegues repetibles. Estas prácticas favorecen ciclos de entrega más rápidos y una administración más eficiente de los entornos tecnológicos \parencite{GoogleCloud2026}.

En consecuencia, la automatización permite transformar procesos tradicionalmente manuales en procesos controlados y reproducibles, contribuyendo a incrementar la eficiencia operativa, mejorar la calidad de los despliegues y facilitar la administración de entornos cloud modernos \parencite{RedHat2026, GoogleCloud2026}.

\subsubsection{Integración entre automatización y monitoreo}

El monitoreo y la automatización constituyen prácticas complementarias dentro de la administración de infraestructura y operaciones en entornos cloud. Mientras que el monitoreo se orienta a recopilar, procesar y analizar información sobre el estado y comportamiento de los sistemas, la automatización utiliza dicha información como entrada para ejecutar acciones de manera automática sobre la infraestructura o los servicios administrados.

En este contexto, los datos obtenidos mediante métricas, registros de eventos y alertas generadas por los mecanismos de monitoreo pueden utilizarse para activar procesos automatizados orientados a responder ante determinadas condiciones operativas. Entre estas acciones pueden encontrarse el aprovisionamiento dinámico de recursos, el escalado automático de infraestructura, la recuperación ante fallos, la ejecución de despliegues o la optimización del consumo de recursos \parencite{Ward2014, RedHat2026}.

La integración entre monitoreo y automatización permite reducir la intervención manual sobre tareas operativas repetitivas y favorecer una administración más dinámica de los entornos tecnológicos. Asimismo, esta relación contribuye a disminuir tiempos de respuesta frente a incidentes y facilita la adaptación del sistema ante cambios en la demanda o en las condiciones de operación \parencite{GoogleCloud2026}.

En consecuencia, el monitoreo proporciona la visibilidad necesaria para comprender el comportamiento del entorno, mientras que la automatización transforma dicha información en acciones ejecutables que permitan mantener el funcionamiento eficiente y continuo de la infraestructura.

\section{Estado del Arte}

\subsection{Soluciones Comerciales Vigentes}

El mercado de herramientas de gestión y optimización cloud ha experimentado un crecimiento significativo en los últimos años. Entre las soluciones comerciales más destacadas se encuentran:

\subsubsection{Amazon CloudWatch y AWS Cost Explorer}

Amazon CloudWatch es un servicio de monitoreo y observabilidad proporcionado por Amazon Web Services que permite recopilar y supervisar métricas, registros (logs) y eventos generados por aplicaciones, servicios e infraestructura desplegada en la nube. La plataforma ofrece funcionalidades de visualización mediante paneles de control, generación de alertas y análisis del comportamiento de los recursos con el objetivo de facilitar la supervisión operativa y la detección de incidentes \parencite{AWS2026}.

Complementariamente, AWS Cost Explorer es una herramienta orientada al análisis y administración de costos que permite visualizar el consumo de servicios cloud, identificar tendencias de gasto, generar reportes y realizar estimaciones de costos futuros. Este servicio proporciona recomendaciones relacionadas con la optimización del consumo dentro del ecosistema AWS \parencite{AWS2026}.

AWS Trusted Advisor es una herramienta que analiza la configuración y utilización de los recursos desplegados en AWS con el objetivo de proporcionar recomendaciones relacionadas con optimización de costos, rendimiento, seguridad, tolerancia a fallos y límites de servicio. Para ello, evalúa la infraestructura mediante un conjunto de verificaciones predefinidas (checks) y genera sugerencias orientadas a identificar recursos subutilizados, configuraciones mejorables o posibles riesgos operativos \parencite{AWS2026}.

En conjunto, Amazon CloudWatch, AWS Cost Explorer y AWS Trusted Advisor proporcionan capacidades de monitoreo, análisis de costos y generación de recomendaciones para la administración de recursos dentro del ecosistema AWS. Mientras CloudWatch permite supervisar métricas, registros y eventos operativos, Cost Explorer facilita el análisis del gasto asociado al consumo de servicios cloud y Trusted Advisor complementa estas funcionalidades mediante verificaciones y recomendaciones predefinidas sobre la infraestructura. Sin embargo, las recomendaciones disponibles dependen de los controles implementados por la plataforma y se encuentran asociadas específicamente a los servicios y recursos soportados por AWS \parencite{AWS2026}.

\subsubsection{Google Cloud Operations Suite}

Google Cloud Operations Suite es el conjunto de herramientas de monitoreo, registro y diagnóstico de Google Cloud destinado a supervisar aplicaciones, servicios e infraestructura desplegados en la nube. La plataforma integra funcionalidades de monitoreo de métricas, gestión de registros (logging), trazabilidad distribuida (tracing), generación de alertas y análisis del rendimiento de sistemas, permitiendo obtener visibilidad sobre el comportamiento operativo de los recursos y facilitar la identificación de incidentes \parencite{GoogleCloud2026}.

La suite proporciona capacidades de observabilidad orientadas a recopilar y correlacionar información proveniente de distintos componentes de la infraestructura, contribuyendo al análisis del estado de los sistemas y al diagnóstico de problemas de rendimiento y disponibilidad. Entre sus funcionalidades se incluyen paneles de visualización, monitoreo de aplicaciones, análisis de eventos y herramientas para la administración de alertas \parencite{GoogleCloud2026}.

Google Cloud Active Assist es un conjunto de herramientas de análisis y recomendación que utiliza datos operativos y técnicas de aprendizaje automático para ayudar a los usuarios a optimizar recursos, mejorar el rendimiento, reforzar la seguridad y administrar costos dentro de Google Cloud. La plataforma analiza continuamente el comportamiento de la infraestructura y genera recomendaciones sobre distintos servicios, incluyendo máquinas virtuales, bases de datos, cuotas de recursos e identidades y accesos, con el objetivo de favorecer una utilización más eficiente de los recursos disponibles. Entre sus funcionalidades se encuentran la identificación de recursos ociosos o sobredimensionados, la detección de configuraciones que podrían afectar el rendimiento de los sistemas y la generación de sugerencias orientadas a la reducción de costos operativos. Estas recomendaciones son presentadas al usuario mediante interfaces integradas dentro de Google Cloud, facilitando la toma de decisiones sobre la infraestructura desplegada \parencite{GoogleCloud2026}.

En conjunto, Google Cloud Operations Suite y Active Assist proporcionan capacidades avanzadas de observabilidad, monitoreo y análisis para recursos desplegados en Google Cloud. Mientras Operations Suite se orienta a la recopilación y visualización de métricas, registros y trazas, Active Assist complementa estas funcionalidades mediante un conjunto de recomendadores (recommenders) especializados que analizan el comportamiento de determinados servicios y generan sugerencias de optimización relacionadas con costos, rendimiento, utilización de recursos y configuración de infraestructura. Sin embargo, las recomendaciones disponibles dependen de los tipos de recursos y recomendadores implementados por la plataforma, y requieren la existencia de un historial de uso suficiente para que el sistema pueda analizar patrones de comportamiento y generar insights relevantes \parencite{GoogleCloud2026}.

\subsubsection{Azure Monitor y Azure Advisor}

Azure Monitor es la plataforma de monitoreo y observabilidad de Microsoft Azure, diseñada para recopilar, almacenar, analizar y visualizar información proveniente de aplicaciones, servicios e infraestructura desplegados en la nube. La herramienta permite supervisar el comportamiento de recursos tanto a nivel de infraestructura como de aplicaciones, proporcionando información sobre rendimiento, disponibilidad y utilización de recursos mediante la recopilación continua de métricas, registros y eventos operativos \parencite{Microsoft2026}.

La plataforma integra diferentes mecanismos de observabilidad que permiten centralizar información generada por máquinas virtuales, bases de datos, redes, contenedores y aplicaciones. Entre sus funcionalidades se encuentran el monitoreo en tiempo real, la generación de alertas automáticas, la visualización mediante paneles de control personalizables y el análisis histórico de métricas. Además, Azure Monitor incorpora herramientas para la gestión de registros (Log Analytics), trazabilidad distribuida de aplicaciones y análisis de eventos, facilitando el diagnóstico de incidentes y la identificación de problemas de rendimiento dentro de arquitecturas distribuidas \parencite{Microsoft2026}.

Azure Advisor es un servicio de recomendaciones que analiza la configuración y utilización de los recursos desplegados en Microsoft Azure con el objetivo de identificar oportunidades de mejora dentro de la infraestructura. A partir de la evaluación continua de los servicios utilizados, la herramienta genera recomendaciones personalizadas orientadas a optimizar costos, incrementar el rendimiento de las aplicaciones, mejorar la disponibilidad de los sistemas, reforzar la seguridad y promover buenas prácticas operativas \parencite{Microsoft2026}.

Las recomendaciones son presentadas al usuario a través del portal de Azure y se encuentran organizadas en distintas categorías, entre las que se incluyen confiabilidad, seguridad, rendimiento, excelencia operativa y optimización de costos. Estas sugerencias pueden abarcar acciones como la identificación de recursos subutilizados, el ajuste de configuraciones, la adopción de mecanismos de alta disponibilidad o la implementación de medidas de seguridad recomendadas por Microsoft. De esta manera, Azure Advisor actúa como un mecanismo de asistencia para la toma de decisiones relacionadas con la administración de infraestructura cloud \parencite{Microsoft2026}.

En conjunto, Azure Monitor y Azure Advisor proporcionan capacidades integradas de observabilidad, monitoreo y generación de recomendaciones para la administración de recursos dentro del ecosistema Microsoft Azure. Mientras Azure Monitor se orienta a la recopilación, almacenamiento y análisis de información operativa proveniente de la infraestructura y las aplicaciones, Azure Advisor utiliza dicha información junto con análisis propios de la plataforma para generar recomendaciones orientadas a mejorar el rendimiento, la seguridad, la disponibilidad y la eficiencia económica de los recursos desplegados. Sin embargo, las recomendaciones disponibles dependen de los mecanismos de análisis implementados por Microsoft y se encuentran asociadas específicamente a los servicios soportados dentro del entorno Azure \parencite{Microsoft2026}.

\subsubsection{Huawei Cloud Eye y Huawei Cloud Cost Center}

Huawei Cloud Eye es el servicio de monitoreo de Huawei Cloud destinado a supervisar el estado y comportamiento de los recursos desplegados dentro de la plataforma. La herramienta permite recopilar métricas de rendimiento provenientes de distintos servicios cloud, incluyendo instancias de cómputo, almacenamiento, redes y bases de datos, proporcionando visibilidad sobre la utilización y disponibilidad de la infraestructura \parencite{HuaweiCloud2026}.

Entre sus principales funcionalidades se encuentran la recolección continua de métricas, la visualización de información mediante gráficos y paneales de control, la configuración de alarmas y notificaciones automáticas, y el análisis histórico del comportamiento de los recursos. Estas capacidades permiten a los administradores identificar problemas operativos, supervisar el consumo de recursos y realizar tareas de seguimiento sobre el rendimiento de los servicios desplegados en la nube \parencite{HuaweiCloud2026}.

A su vez, Cloud Eye se integra con otros servicios del ecosistema Huawei Cloud para centralizar la supervisión de la infraestructura y facilitar la administración operativa del entorno. Es decir, la plataforma proporciona información relevante para la detección de incidentes y el monitoreo continuo de sistemas cloud, constituyendo uno de los principales mecanismos de observabilidad disponibles dentro del proveedor. Huawei Cloud Cost Center es una herramienta orientada a la gestión financiera de recursos cloud que permite visualizar, analizar y controlar los costos asociados al consumo de servicios dentro de Huawei Cloud. La plataforma proporciona información sobre gastos históricos, tendencias de consumo y distribución de costos entre distintos recursos y servicios utilizados por la organización \parencite{HuaweiCloud2026}.

Entre sus funcionalidades se incluyen el seguimiento del gasto cloud, la generación de reportes de costos, la definición de presupuestos y el monitoreo del consumo de recursos desde una perspectiva financiera. Estas capacidades permiten obtener mayor visibilidad sobre el impacto económico de la infraestructura desplegada y facilitan la adopción de prácticas de control y optimización del gasto tecnológico \parencite{HuaweiCloud2026}.

En conjunto, Huawei Cloud Eye y Huawei Cloud Cost Center proporcionan capacidades de monitoreo operativo y gestión financiera para la administración de recursos dentro del ecosistema Huawei Cloud. Mientras Cloud Eye se enfoca en la recopilación y visualización de métricas relacionadas con el rendimiento y disponibilidad de la infraestructura, Cost Center permite analizar los costos asociados a la utilización de los servicios cloud. Sin embargo, las funcionalidades ofrecidas por estas herramientas se encuentran principalmente orientadas al monitoreo y análisis de información, sin incorporar un mecanismo integral de detección de anomalías basado en aprendizaje automático no supervisado que permita identificar patrones inusuales de utilización de recursos y generar recomendaciones específicas de optimización sobre la infraestructura desplegada \parencite{HuaweiCloud2026}.

\subsubsection{VMware CloudHealth}

VMware CloudHealth es una plataforma de gestión y optimización de entornos cloud diseñada para proporcionar visibilidad sobre el consumo de recursos, los costos operativos y el cumplimiento de políticas dentro de infraestructuras desplegadas en múltiples proveedores cloud. La herramienta recopila información proveniente de servicios de computación, almacenamiento, redes y otros recursos cloud con el objetivo de centralizar su análisis y facilitar la administración de entornos complejos \parencite{VMware2023}.

Entre sus principales funcionalidades se encuentran el monitoreo del consumo de recursos, el análisis de costos cloud, la asignación de gastos entre equipos o proyectos y la generación de reportes orientados a la adopción de prácticas FinOps. La plataforma permite establecer políticas de gobernanza y supervisar el cumplimiento de configuraciones definidas por la organización, proporcionando una visión integral tanto desde una perspectiva técnica como financiera \parencite{VMware2023}.

CloudHealth incorpora capacidades de optimización mediante la generación de recomendaciones destinadas a mejorar la eficiencia en la utilización de recursos y reducir costos operativos. Estas recomendaciones se basan en el análisis de métricas históricas de consumo y en la identificación de patrones asociados a recursos infrautilizados, instancias sobredimensionadas, almacenamiento sin uso y otras oportunidades de optimización. Por esto, la plataforma asiste a los administradores en la toma de decisiones relacionadas con la gestión y el dimensionamiento de la infraestructura cloud \parencite{VMware2023}.

En conjunto, las funcionalidades de monitoreo, análisis financiero y optimización de VMware CloudHealth proporcionan una solución orientada a la gestión eficiente de recursos cloud dentro de entornos empresariales. Sin embargo, sus capacidades de recomendación se encuentran principalmente enfocadas en la optimización financiera y la gobernanza de infraestructura, mientras que el análisis realizado por la plataforma se basa en mecanismos y criterios definidos por el proveedor. Su adopción se encuentra orientada principalmente a los ecosistemas cloud de mayor presencia en el mercado, como AWS, Microsoft Azure y Google Cloud Platform \parencite{VMware2023}.

\subsubsection{Datadog}

Datadog es una plataforma de monitoreo y seguridad basada en la nube que proporciona observabilidad integral para aplicaciones, infraestructura y servicios en entornos cloud modernos. La plataforma integra capacidades de monitoreo de infraestructura, aplicaciones, logs, seguridad y experiencia digital en una solución unificada, permitiendo a las organizaciones obtener visibilidad completa sobre sus stacks tecnológicos a cualquier escala \parencite{Datadog2026}.

Entre sus principales funcionalidades se encuentran el monitoreo de infraestructura con visualización detallada de hosts, contenedores y recursos cloud; Application Performance Monitoring (APM) con trazabilidad distribuida y perfiles de código; gestión de logs con análisis y correlación automática; monitoreo de experiencia digital mediante Real User Monitoring (RUM) y Synthetic Monitoring; y capacidades de seguridad cloud que incluyen Cloud Security Posture Management (CSPM) y detección de amenazas \parencite{Datadog2026}.

Datadog incorpora funcionalidades avanzadas de gestión de costos cloud mediante Cloud Cost Management, que permite visualizar y optimizar el gasto en múltiples proveedores cloud incluyendo AWS, Azure y Google Cloud. La plataforma proporciona recomendaciones de optimización de costos basadas en el análisis de patrones de uso, identificación de recursos infrautilizados y sugerencias de dimensionamiento adecuado de instancias \parencite{Datadog2026}.

La plataforma destaca por su integración nativa con más de 800 tecnologías y servicios, incluyendo Kubernetes, Docker, AWS, Azure, Google Cloud, bases de datos, sistemas de mensajería y herramientas de CI/CD. Además, incorpora capacidades de inteligencia artificial mediante Bits AI, que proporciona asistentes impulsados por IA para investigación, análisis y remediación automática de incidentes \parencite{Datadog2026}.

En conjunto, Datadog proporciona una solución integral de observabilidad que combina monitoreo de infraestructura, aplicaciones, logs, seguridad y experiencia digital en una plataforma unificada. Su enfoque multi-cloud y las capacidades de correlación automática entre diferentes fuentes de telemetría permiten obtener una visión completa del comportamiento de sistemas distribuidos. Sin embargo, su modelo de pricing basado en volumen de datos ingeridos puede resultar costoso para organizaciones con grandes volúmenes de telemetría, y las recomendaciones de optimización se basan principalmente en reglas predefinidas y análisis estadísticos \parencite{Datadog2026}.

\subsubsection{Dynatrace}

Dynatrace es una plataforma de observabilidad impulsada por inteligencia artificial diseñada para proporcionar visibilidad completa sobre aplicaciones, infraestructura cloud y experiencia digital. La plataforma utiliza un enfoque automático de instrumentación mediante su agente OneAgent, que despliega automáticamente la instrumentación para aplicaciones, infraestructura y servicios sin requerir configuración manual extensiva \parencite{Dynatrace2026}.

Entre sus principales funcionalidades se encuentran el monitoreo de aplicaciones con trazabilidad distribuida automática y análisis de rendimiento de código; observabilidad de infraestructura para entornos multi-cloud e híbridos; análisis de logs con capacidades de búsqueda y correlación; monitoreo de experiencia digital mediante Real User Monitoring y Synthetic Monitoring; y observabilidad de inteligencia artificial para aplicaciones que utilizan LLMs y agentes de IA \parencite{Dynatrace2026}.

Dynatrace incorpora Davis AI, un motor de inteligencia artificial determinista que proporciona detección automática de anomalías, análisis de causa raíz y predicción de problemas. A diferencia de enfoques basados en machine learning probabilístico, Davis AI utiliza modelos causales deterministas que permiten identificar con precisión la causa raíz de problemas y generar recomendaciones accionables sin generar falsos positivos \parencite{Dynatrace2026}.

La plataforma proporciona capacidades avanzadas de automatización mediante Dynatrace Intelligence, que permite predecir problemas, detectar anomalías y coordinar acciones automatizadas a través de plataformas cloud, herramientas de desarrollo y sistemas de gestión de servicios de TI. Además, incorpora Grail, un data lakehouse que permite almacenar y analizar datos de telemetría en contexto, facilitando consultas complejas y análisis avanzados \parencite{Dynatrace2026}.

En conjunto, Dynatrace proporciona una solución de observabilidad con enfoque en automatización e inteligencia artificial determinista. Su capacidad de instrumentación automática reduce la complejidad de configuración, mientras que Davis AI proporciona análisis causal preciso para identificación de problemas. Sin embargo, su enfoque propietario y el modelo de pricing pueden representar una inversión significativa para organizaciones, y las capacidades de optimización de costos cloud son menos desarrolladas comparadas con soluciones especializadas en FinOps \parencite{Dynatrace2026}.

\subsubsection{New Relic}

New Relic es una plataforma de observabilidad inteligente que proporciona visibilidad completa sobre el stack tecnológico, desde infraestructura hasta experiencia del usuario. La plataforma integra monitoreo de aplicaciones, infraestructura, logs, experiencia digital y seguridad en una solución unificada, permitiendo correlacionar datos de diferentes fuentes para obtener una visión integral del comportamiento de los sistemas \parencite{NewRelic2026}.

Entre sus principales funcionalidades se encuentran Application Performance Monitoring (APM) con trazabilidad distribuida y análisis de transacciones; monitoreo de infraestructura para entornos cloud, Kubernetes y sistemas híbridos; gestión de logs con análisis inteligente y correlación automática; monitoreo de experiencia digital mediante Browser Monitoring, Mobile Monitoring y Synthetic Monitoring; y capacidades de seguridad con gestión de vulnerabilidades \parencite{NewRelic2026}.

New Relic incorpora capacidades avanzadas de inteligencia artificial mediante New Relic AI y AIOps, que proporcionan detección de anomalías, predicción de incidentes y automatización de respuestas. La plataforma incluye SRE Agent, un agente impulsado por IA que automatiza la investigación y remediación de incidentes, reduciendo el tiempo de resolución y la carga operativa sobre los equipos de ingeniería \parencite{NewRelic2026}.

La plataforma proporciona Cloud Cost Intelligence, una funcionalidad que permite visualizar y optimizar costos en entornos multi-cloud y Kubernetes. Esta herramienta proporciona visibilidad en tiempo real sobre el gasto cloud, identifica oportunidades de optimización y permite correlacionar costos con el rendimiento y utilización de recursos. Además, New Relic soporta OpenTelemetry como estándar abierto para la recopilación de telemetría \parencite{NewRelic2026}.

En conjunto, New Relic proporciona una plataforma de observabilidad con enfoque en inteligencia artificial y automatización. Su modelo de pricing basado en datos ingeridos ofrece flexibilidad para organizaciones de diferentes tamaños, y las capacidades de correlación automática facilitan el análisis de sistemas distribuidos. Sin embargo, la profundidad de las recomendaciones de optimización de costos puede ser limitada comparada con soluciones especializadas, y la configuración inicial puede requerir esfuerzo significativo para obtener observabilidad completa \parencite{NewRelic2026}.

\subsection{Soluciones de Código Abierto}

Además de las soluciones comerciales, existen herramientas de código abierto que ofrecen capacidades de monitoreo y observabilidad para entornos cloud:

\subsubsection{Prometheus}

Prometheus es un sistema de monitoreo y alertamiento de código abierto que recopila métricas de aplicaciones y servicios mediante un modelo de extracción (pull model). La herramienta utiliza un modelo de datos multidimensional con identificadores de métricas y etiquetas, permitiendo consultas flexibles mediante su lenguaje PromQL. Prometheus es ampliamente adoptado en entornos Kubernetes y proporciona capacidades de almacenamiento de series temporales, generación de alertas y visualización mediante integración con herramientas como Grafana \parencite{Prometheus2026}.

\subsubsection{Grafana}

Grafana es una plataforma de visualización y análisis de código abierto que permite crear paneles de control interactivos para métricas provenientes de múltiples fuentes de datos. La herramienta soporta integración con Prometheus, InfluxDB, Elasticsearch y otros sistemas de almacenamiento de métricas, proporcionando capacidades de visualización, alertamiento y análisis de datos operativos \parencite{Grafana2026}.

\subsubsection{OpenTelemetry}

OpenTelemetry es un proyecto de código abierto que proporciona un conjunto de APIs, SDKs y herramientas para la instrumentación, generación, recopilación y exportación de datos de telemetría (métricas, registros y trazas). El proyecto busca proporcionar un estándar vendor-neutral para la observabilidad, permitiendo a las organizaciones recopilar datos de telemetría independientemente del proveedor cloud o del backend de análisis utilizado \parencite{OpenTelemetry2026}.

\subsection{Antecedentes de Investigación}

Además de las soluciones comerciales y de código abierto, existen antecedentes académicos que abordan la detección de anomalías en infraestructura cloud mediante técnicas de aprendizaje automático, los cuales resultan relevantes para contextualizar la propuesta de esta tesis.

\subsubsection{Zou et al. (2022)}

Zou et al. (2022) desarrollaron un sistema de monitoreo de anomalías para contenedores Docker en entornos de computación en la nube, basado en una versión optimizada del algoritmo Isolation Forest \parencite{Zou2022}. El sistema analiza métricas multidimensionales asociadas a la utilización de recursos, incluyendo CPU, memoria, disco y red, con el objetivo de identificar comportamientos anómalos durante la ejecución de los contenedores. Como parte de su propuesta, los autores modifican el mecanismo de selección de características de Isolation Forest mediante la asignación de diferentes pesos a las métricas de acuerdo con su relevancia para la carga de trabajo analizada. El sistema fue evaluado tanto en entornos simulados como en entornos cloud reales, demostrando su capacidad para detectar diferentes tipos de anomalías asociadas a los recursos computacionales.

Este trabajo presenta una relación directa con InfraEye, debido a que ambas propuestas utilizan técnicas de aprendizaje no supervisado para analizar métricas de utilización de recursos y detectar comportamientos atípicos mediante Isolation Forest. Sin embargo, mientras Zou et al. se enfocan en la detección de anomalías operativas en contenedores, InfraEye orienta el análisis hacia la identificación de comportamientos sostenidos de utilización ineficiente de recursos cloud, con el propósito de determinar posibles situaciones de sobredimensionamiento o subutilización. A partir de esta detección, InfraEye incorpora además un mecanismo de recomendación de recursos y automatización mediante Infrastructure as Code, aspectos que no forman parte del objetivo principal del trabajo de Zou et al.

A diferencia de la propuesta de Zou et al., que utiliza las métricas de utilización como características individuales del modelo, InfraEye incorpora características estadísticas calculadas sobre ventanas temporales de las métricas recolectadas. Esta estrategia permite caracterizar el comportamiento de utilización durante períodos prolongados y resulta consistente con el objetivo de identificar situaciones persistentes de subutilización o sobredimensionamiento, en lugar de centrarse exclusivamente en eventos anómalos puntuales.

\subsubsection{Nätterdal y Olausson (2024)}

Nätterdal y Olausson (2024) desarrollaron un sistema de detección de anomalías para entornos cloud basado en Isolation Forest, aplicado sobre métricas de monitoreo obtenidas mediante Amazon CloudWatch \parencite{Natterdal2024}. El trabajo aborda específicamente las dificultades asociadas con la utilización de Isolation Forest sobre datos de streaming, considerando fenómenos como cambios estacionales y concept drift. Para ello, los autores desarrollaron un modelo basado en Extended Isolation Forest, complementado con técnicas de preprocesamiento, agrupamiento mediante K-Means y detección de cambios mediante ADWIN. La propuesta fue evaluada sobre métricas reales de Amazon CloudWatch, recolectadas cada cinco minutos durante un período de aproximadamente tres meses, y posteriormente desplegada en AWS para realizar detección sobre datos en tiempo real.

Este trabajo presenta una relación directa con InfraEye debido a que ambas propuestas emplean técnicas basadas en Isolation Forest para analizar métricas de infraestructura cloud y detectar comportamientos que se apartan de los patrones considerados normales. Sin embargo, mientras Nätterdal y Olausson se enfocan en la detección de anomalías operativas sobre flujos de datos en tiempo real y en la adaptación del modelo frente a cambios en la distribución de los datos, InfraEye utiliza características estadísticas calculadas sobre ventanas temporales con el objetivo de identificar comportamientos sostenidos de utilización ineficiente de recursos. A partir de esta identificación, InfraEye extiende el proceso hacia la recomendación de recursos cloud y la generación de modificaciones de infraestructura mediante Terraform, orientando el análisis hacia la optimización del dimensionamiento y del uso de los recursos.

\subsubsection{Nawrocki y Sus (2022)}

Nawrocki y Sus (2022) desarrollaron un sistema para la planificación a largo plazo de recursos cloud basado en la detección de anomalías y técnicas de aprendizaje automático \parencite{Nawrocki2022}. La propuesta analiza métricas de utilización de CPU, memoria, red y disco de máquinas virtuales, utilizando un mecanismo híbrido de detección denominado Weighted Hybrid Algorithm (WHA), diseñado para identificar anomalías puntuales y recurrentes y adaptarse a cambios en los patrones de utilización. Los datos procesados mediante este mecanismo son posteriormente utilizados para predecir la demanda futura y ajustar la asignación de recursos, con el objetivo de reducir situaciones de sobredimensionamiento y los costos asociados. La solución fue evaluada utilizando datos reales correspondientes a más de 1700 máquinas virtuales.

Este trabajo presenta una relación directa con InfraEye debido a que ambas propuestas utilizan métricas de utilización de recursos cloud y detección de anomalías como parte del proceso de optimización del dimensionamiento. Sin embargo, se identifican diferencias significativas tanto en el enfoque de detección como en el uso posterior de los resultados. En primer lugar, Nawrocki y Sus emplean un mecanismo híbrido de detección (WHA) que combina múltiples algoritmos ponderados para identificar tanto anomalías puntuales como recurrentes, con capacidad de adaptación ante cambios en los patrones de utilización. InfraEye, en cambio, utiliza Isolation Forest como algoritmo único, aplicado sobre características estadísticas calculadas sobre ventanas temporales, orientado a detectar comportamientos sostenidos de ineficiencia más que eventos anómalos aislados. En segundo lugar, la propuesta de Nawrocki y Sus utiliza las anomalías detectadas como insumo para un modelo de predicción de demanda futura de recursos, con el objetivo de anticipar necesidades de capacidad y ajustar la asignación de recursos de manera anticipada. InfraEye no incorpora predicción: el análisis se circunscribe a la identificación de ineficiencias presentes en los datos históricos, a partir de las cuales se generan recomendaciones concretas de redimensionamiento orientadas a corregir situaciones actuales de sobredimensionamiento o subutilización. En tercer lugar, la evaluación de Nawrocki y Sus fue realizada sobre un conjunto de más de 1700 máquinas virtuales, lo cual proporciona evidencia sobre la escalabilidad del enfoque híbrido propuesto, mientras que InfraEye fue validado en un entorno controlado con un alcance más acotado. Finalmente, InfraEye incorpora un módulo de integración con Terraform que permite automatizar la aplicación de las modificaciones de infraestructura recomendadas, cerrando el ciclo entre detección y acción, aspecto que no forma parte de la propuesta de Nawrocki y Sus.

\subsection{Análisis Comparativo}

El análisis realizado muestra que las principales soluciones comerciales de gestión cloud, como AWS CloudWatch y Trusted Advisor, Google Cloud Operations Suite y Active Assist, Azure Monitor y Advisor, Huawei Cloud Eye con Cost Center, Datadog, Dynatrace y New Relic ofrecen capacidades avanzadas de monitoreo, observabilidad y optimización de recursos. Estas plataformas permiten supervisar métricas operativas, visualizar el consumo de infraestructura y generar recomendaciones orientadas a reducir costos o mejorar el rendimiento de los sistemas.

No obstante, también se identifican limitaciones comunes que restringen su alcance frente a las necesidades planteadas en el presente proyecto. En primer lugar, existe una fuerte dependencia del proveedor cloud en las herramientas nativas, ya que se encuentran diseñadas para operar principalmente dentro de sus propios ecosistemas, dificultando la portabilidad y la adaptación a otros entornos. Esta situación resulta especialmente relevante en el caso de Huawei Cloud, donde la oferta de herramientas avanzadas de optimización es más limitada que en AWS, Azure o Google Cloud.

En segundo lugar, las recomendaciones generadas por estas plataformas se basan mayormente en reglas, verificaciones o recommenders predefinidos por el proveedor. Si bien algunos servicios, como Active Assist y Dynatrace Davis AI, incorporan mecanismos de análisis más sofisticados, el usuario no dispone de control ni visibilidad sobre el modelo utilizado ni puede adaptar fácilmente el comportamiento de detección a patrones específicos de su organización. En contraste, el presente proyecto propone un enfoque explícito y documentado basado en Isolation Forest, un algoritmo de aprendizaje automático no supervisado orientado a identificar anomalías directamente a partir de las métricas operativas recolectadas.

Asimismo, las soluciones comerciales analizadas presentan una escasa integración con prácticas de Infraestructura como Código (IaC). Aunque permiten detectar oportunidades de optimización, la implementación de los cambios recomendados suele requerir intervención manual o herramientas externas. La propuesta de esta tesis incorpora un módulo integrado con Terraform, permitiendo automatizar el despliegue y la administración de infraestructura desde la misma plataforma.

Otra diferencia importante radica en el enfoque del análisis. Las herramientas existentes se orientan principalmente a un modelo reactivo basado en umbrales y reglas predefinidas, donde los problemas se detectan una vez que ya se manifiestan en la infraestructura o en el consumo. La solución propuesta utiliza aprendizaje automático no supervisado para detectar anomalías e ineficiencias en los datos históricos de utilización que los mecanismos reactivos basados en umbrales y reglas predefinidas no logran identificar, generando recomendaciones sobre comportamientos anómalos que no constituyen incidentes manifiestos pero representan oportunidades de optimización.

En síntesis, el estado del arte revela que tanto el mercado como la investigación académica ofrecen aportes relevantes pero insuficientes para cubrir las necesidades operativas identificadas. Por un lado, el mercado proporciona soluciones con capacidades consolidadas de monitoreo, observabilidad y generación de recomendaciones de optimización, tanto en herramientas nativas de proveedores cloud (Amazon CloudWatch y Trusted Advisor, Google Cloud Operations Suite y Active Assist, Azure Monitor y Advisor, Huawei Cloud Eye y Cost Center) como en plataformas multi-cloud comerciales (Datadog, Dynatrace, New Relic, VMware CloudHealth) y herramientas de código abierto (Prometheus, Grafana, OpenTelemetry). Por otro lado, la investigación académica ha abordado la detección de anomalías en infraestructura cloud mediante Isolation Forest: Zou et al. (2022) propusieron una versión optimizada para contenedores Docker con selección ponderada de características; Nätterdal y Olausson (2024) desarrollaron un modelo basado en Extended Isolation Forest con técnicas de adaptación ante concept drift sobre datos de streaming; y Nawrocki y Sus (2022) combinaron detección de anomalías mediante un mecanismo híbrido (WHA) con predicción de demanda futura para planificación de recursos. Sin embargo, persisten brechas tanto en las soluciones de mercado como en los antecedentes académicos. Respecto del mercado, se identifican tres limitaciones fundamentales: la dependencia del proveedor cloud, que restringe la portabilidad y resulta especialmente crítica en Huawei Cloud; las recomendaciones basadas en reglas y verificadores predefinidos, sin control ni transparencia sobre los modelos subyacentes; y la escasa integración con Infraestructura como Código, que requiere intervención manual para implementar cambios. Respecto de la investigación, los antecedentes académicos presentan limitaciones complementarias: se orientan a la detección de anomalías operativas puntuales o en tiempo real (Zou et al., Nätterdal y Olausson) o a la predicción de demanda futura (Nawrocki y Sus), pero no abordan la identificación de comportamientos sostenidos de ineficiencia seguida de la generación de recomendaciones de dimensionamiento y su automatización mediante IaC. INFRAEYE se propone cubrir estas brechas de manera integral mediante tres diferenciadores concretos: (a) la aplicación de Isolation Forest como algoritmo de aprendizaje automático no supervisado explícito y documentado, que permite detectar anomalías en el comportamiento histórico de las métricas operativas que los enfoques basados en reglas predefinidas no logran capturar, sin depender de verificadores del proveedor ni requerir datos etiquetados, y que a diferencia de los antecedentes académicos opera sobre características estadísticas calculadas sobre ventanas temporales para identificar comportamientos sostenidos de ineficiencia; (b) la integración nativa con Terraform, que habilita la automatización del despliegue y la administración de infraestructura desde la propia plataforma, cerrando el ciclo entre detección y acción, aspecto ausente tanto en las soluciones comerciales como en los antecedentes académicos; y (c) la capacidad de identificar ineficiencias en la infraestructura existente a partir del análisis de datos históricos de utilización, generando recomendaciones accionables de redimensionamiento sobre recursos que presentan comportamientos anómalos respecto del patrón operacional de referencia, sin recurrir a predicción de demanda futura. De este modo, INFRAEYE complementa tanto las capacidades del mercado como los aportes de la investigación previa al articular monitoreo continuo, detección inteligente de anomalías basada en datos históricos, recomendaciones adaptables al contexto operacional y automatización de infraestructura en una solución unificada orientada al ecosistema Huawei Cloud.
